\subsection{Abordagem Metodológica}
O desenho da arquitetura do sistema seguiu a metodologia de modelação dimensional orientada aos processos de negócio, proposta por Ralph Kimball. Esta abordagem permitiu identificar os processos fundamentais, definindo para cada um o grão (nível de detalhe), as dimensões (contexto analítico) e as medidas (factos numéricos).

Com base nos requisitos levantados, foram modelados quatro processos de negócio distintos, conforme descrito abaixo.

\subsubsection{Ocorrência de Incêndio}
Este processo centraliza a informação relativa aos eventos de fogo.
\begin{itemize}
	\item \textbf{Grão:} Um registo por incêndio, identificado univocamente pelo \texttt{incident\_id}.
	\item \textbf{Dimensões:} Tempo (início/fim), Localização (Distrito, Concelho, Freguesia), Tipo de Incêndio e possível causa.
	\item \textbf{Medidas:} Área ardida (ha), duração (horas) e métricas derivadas (ex: taxa de expansão).
\end{itemize}

\subsubsection{Meteorologia Diária}
Permite a análise das condições ambientais que propiciam ou agravam as ocorrências.
\begin{itemize}
	\item \textbf{Grão:} Registo diário das condições meteorológicas num determinado local geográfico afetado.
	\item \textbf{Dimensões:} Tempo (Dia) e Localização (coordenadas aproximadas ou freguesia).
	\item \textbf{Medidas:} Temperatura (média e máxima), humidade relativa, velocidade do vento e precipitação acumulada.
\end{itemize}

\subsubsection{Cobertura Mediática}
Analisa a repercussão pública e o impacto mediático dos eventos.
\begin{itemize}
	\item \textbf{Grão:} Um artigo noticioso individual, identificado univocamente por um id.
	\item \textbf{Dimensões:} Tempo (Data de publicação), Fonte (Entidade emissora/Jornal) e Sentimento (Classificação da tonalidade da notícia).
	\item \textbf{Medidas:} \textit{Risk Score} (pontuação de risco associada), Polaridade (valor numérico do sentimento) e volume de notícias (contagem).
\end{itemize}

\subsubsection{Monitorização de Vegetação}
Avalia o impacto ecológico e o estado do combustível florestal.
\begin{itemize}
	\item \textbf{Grão:} Medições de índices de vegetação agregadas por perímetro de incêndio e janela temporal.
	\item \textbf{Dimensões:} Tempo (Janela pré/pós incêndio), Localização.
	\item \textbf{Medidas:} NDVI médio pré-incêndio (combustível disponível), NDVI médio pós-incêndio e variação percentual (severidade do dano).
\end{itemize}